\section{第\ref{sec:muscles}章　筋肉への出力}

出力の仕組み、すなわち運動単位、シナプスの安全率、大きさの順の動員、筋肉の対、伸張ループ、リズム発生器は直接測定されている。遅延ループの安定条件は定理である。体の内部モデルは十分な根拠のある仮説である。図の数値は本書のモデルによる計算である。

{\footnotesize
\setlength{\tabcolsep}{3pt}\renewcommand{\arraystretch}{1.15}
\begin{longtable}{>{\raggedright}p{0.5cm}>{\raggedright}p{4.0cm}>{\raggedright}p{5.6cm}>{\raggedright}p{2.3cm}>{\raggedright\arraybackslash}p{1.7cm}}
\toprule
No. & 主張 & 実験と測定されたもの & 出典 & 状態 \\
\midrule\endfirsthead
\toprule
No. & 主張 & 実験と測定されたもの & 出典 & 状態 \\
\midrule\endhead
1 & 運動単位：1個の\nt{ACET}ニューロンが数百から2000本ほどの線維のグループを制御する；細かな調整を要する筋肉では単位が小さい & ヒトの筋肉、運動軸索と筋線維の計数：ニューロンあたりの線維は平均で手の骨間筋で約 $340$、腕橈骨筋で約 $410$、腓腹筋内側頭で約 $1930$。最も小さな単位の正確な数は章では挙げていない。 & \cite{Feinstein1955mu} & 測定済み \\
2 & 筋肉上のシナプスは、線維の発火に必要な量の数倍の伝達物質を放出する & 神経筋シナプスでの測定の総説：安全率（放出された伝達物質と閾値の比）は成体哺乳類で通常 $3$--$5$；ヒトでは安全率は放出量よりも受け手の膜のひだによって保たれている。 & \cite{WoodSlater2001} & 測定済み \\
3 & 単収縮~\eqref{eq:twitch}は $50\ms$ 程度の時間 $T_c$ で立ち上がる & 随意収縮中のヒトの手の単一運動単位、単位のスパイクで加算平均した力：立ち上がり時間 $30$--$100\ms$、$80\,\%$ の単位で $70\ms$ 未満。関数 $(t/T_c)e^{1-t/T_c}$ は運動単位プールのモデルからの近似。 & \cite{MilnerBrown1973rec, Fuglevand1993} & 測定済み \\
4 & 単収縮の和~\eqref{eq:forcesum}：$5$, $15$, $40$ スパイク/s で力は $0.2$--$1.1F_1$、$1.8$--$2.2F_1$、約 $5.4F_1$（図~\ref{fig:tetanus}） & \texttt{models/\allowbreak muscle\_\allowbreak models.py} による計算、$T_c=50\ms$：最後の $0.3\,\text{s}$ で $0.21$--$1.10$、$1.76$--$2.19$、$5.28$--$5.49\,F_1$。線形の加算は単純化：実際の筋肉の飽和はモデルにない。 & \texttt{models/\allowbreak muscle\_\allowbreak models.py} & モデル \\
5 & 単位は大きさの順に動員される；最も弱いものは最も強いものの100分の1 & ネコの前根：反射入力を強めていくと、前根でのスパイクが小さいニューロン、つまり小さいニューロンが最初に発火する。ヒトの手の骨間筋：単位の単収縮の力は $0.1$ から $10$~g で、単位が動員される力にほぼ比例して増える。 & \cite{Henneman1957, MilnerBrown1973rec} & 測定済み \\
6 & 力の刻みは力に比例する、$\Delta F/F\approx q-1$~\eqref{eq:sizeprinciple}：浮動小数点 & 章で力の等比数列から導出。実験と整合する：大きな力では、同じ力の増分に対して新たに動員される単位がずっと少ない。 & 章内の導出；\cite{MilnerBrown1973rec} & 理論 \\
7 & 力のばらつきは力そのものに比例して増える & 母指の筋の随意等尺性収縮：力の標準偏差は平均の力に比例して増える；電気刺激で同じ力を出させるとばらつきは一定。プールのモデル：力の順の動員が線形の増加を生む。 & \cite{Jones2002sdn} & 測定済み \\
8 & 動きのなめらかさは信号依存ノイズの帰結である & モデル：このようなノイズのもとで終点のばらつきを最小にする軌道は、測定された眼と腕の軌道を再現する。この原因を他の原因と区別する直接の実験はない。 & \cite{HarrisWolpert1998} & 仮説 \\
9 & 筋肉の対の力の差がトルクを、和が剛性を決める~\eqref{eq:antagonists} & 同時収縮によるインピーダンス制御の理論。ヒトの腕：小さな摂動で測った各関節の剛性は、その関節のトルクにほぼ比例する；同時収縮の比率を変えると手先の剛性楕円の形と向きが変わる。 & \cite{Hogan1984, GomiOsu1998stiff} & 測定済み \\
10 & 伸張センサは自分の\nt{ACET}ニューロンを興奮させ、抑制性の介在ニューロンを通じて拮抗筋をオフにする（図~\ref{fig:antagonists}） & ネコの脊髄：伸張センサの線維に興奮させられる単一の介在ニューロンが、拮抗筋の\nt{ACET}ニューロンに単シナプス性の抑制性電位を生む。シナプス遅延 $0.28$--$0.42\ms$。介在ニューロンの伝達物質（グリシン）はこの実験では同定されていない。 & \cite{Jankowska1972Ia} & 測定済み \\
11 & ループの遅延：脊髄 $25$--$30\ms$、皮質経由はその1.5〜3倍、目経由 $100$--$200\ms$ & ヒトの腕、関節への摂動：筋肉の短潜時応答は $20$--$45\ms$、長潜時応答は $45$--$105\ms$、随意応答は $120$--$180\ms$。運動中に見える指の位置をずらすと、修正は平均 $160\ms$ 後。 & \cite{Pruszynski2008rapid, SaundersKnill2003vis} & 測定済み \\
12 & $\dd x/\dd t=-Gx(t-d)$ は $Gd<\pi/2$~\eqref{eq:delaystable}で安定；境界での周波数は $1/(4d)$ & 遅延方程式の根についての定理。\texttt{models/\allowbreak muscle\_\allowbreak models.py} による検証、$d=30\ms$：$3\,\text{s}$ の時点で振幅は $G=40$ で $3\cdot10^{-6}$、$G=50$ で $0.13$、$G=55$ で $38$（単位は毎秒、境界は $52$）。 & \cite{Hayes1950} & 理論 \\
13 & 生理的振戦 $8$--$12$~Hz；伸張ループはその源の一つ & 測定の総説：健康な人には約 $10$~Hz 帯の細かな震えがある；その起源は混合的で、中枢の振動、単位の発火の性質、機械的共振、反射ループの共振があり、条件によって優勢なものが異なる。 & \cite{McAuleyMarsden2000} & 仮説 \\
14 & 脳は体の内部モデルによって指令の結果を予測する & ヒトが物体をつまんで持ち、腕を動かす：慣性、粘性、弾性の負荷のいずれでも、把持力はループの遅れなしに負荷と同時に変わる。つまり負荷は予測されている。総説がこうした実験をまとめる；ニューロン内のモデルそのものの直接観察はない。 & \cite{FlanaganWing1997grip, WolpertGhahramani2000} & 仮説 \\
15 & 歩行、呼吸、咀嚼のリズムは、一定の入力のもとで局所的なネットワークが生む & 実験の総説：摘出した神経系（脊髄、神経節）は、リズミカルな入力も筋肉からのフィードバックもなしにリズミカルな運動パターンを出す。 & \cite{MarderBucher2001} & 測定済み \\
16 & 疲労をともなう相互抑制~\eqref{eq:halfcentre}は周期 $0.84\,\text{s}\approx2\tau_a$ のリズムを生む（図~\ref{fig:halfcentre}） & 定理：相互抑制と適応を持ち、一定の入力のもとで安定状態を持たないネットワークは必ず振動する。\texttt{models/\allowbreak muscle\_\allowbreak models.py} による計算（$I=1$, $\beta=2$, $b=2$, $\tau=20\ms$, $\tau_a=0.4\,\text{s}$）：周期 $0.84\,\text{s}$。実際の発生器がまさにこのような構成であることは示されていない。 & \cite{Matsuoka1985osc} & モデル \\
\bottomrule
\end{longtable}}
